% ------------------------------------------------------------------------------
\newpage
\section{Skalarprodukt und Winkel}

In der Vektorgeometrie benötigen wir eine Möglichkeit, Winkel zwischen Vektoren
zu berechnen. Ausgehend von Betrachtungen am Dreieck und durch einen Vergleich
mit dem Kosinussatz erhalten wir eine Definition für das \textbf{Skalarprodukt}
von Vektoren, welches diese Möglichkeit bietet.

% ------------------------------------------------------------------------------
\subsection{Herleitung aus dem Kosinussatz}
Im folgenden wird gezeigt, wie das Skalarprodukt aus dem Kosinussatz hergeleitet
wird. Dazu wird das folgende Dreieck betrachtet:
\begin{center}
  \begin{tikzpicture}
    \tkzDefPoint(3,2){C}
    \tkzDefPoint(8,1){A}
    \tkzDefPoint(7,4){B}
    \tkzDrawSegment[thick,vector style](C,A)
    \tkzDrawSegment[thick,vector style](C,B)
    \tkzDrawSegment[thick,vector style](A,B)
    \tkzLabelSegment[below left](C,A){$\vv{a}$}
    \tkzLabelSegment[above right](A,B){$\vv{c}$}
    \tkzLabelSegment[above left](C,B){$\vv{b}$}
    \tkzMarkAngle[size=1](A,C,B)
    \tkzLabelAngle[pos=0.7](A,C,B){$\gamma$}

    \tkzDrawPoints(A,B,C)
    \tkzLabelPoint[below right](A){$A$}
    \tkzLabelPoint[above](B){$B$}
    \tkzLabelPoint[below left](C){$C$}
  \end{tikzpicture}
\end{center}
Die Vektoren $\vv{a}$ und $\vv{b}$ werden als Komponenten geschrieben. Der
Vektor $\vv{c}$ kann aus den anderen berechnet werden.
\begin{align*}
  \vv{a} &= \vxy{a_{x},a_{y}} &
  \vv{b} &= \vxy{b_{x},b_{y}} &
  \vv{c} &= \vv{b}-\vv{a} = \vxy{b_{x}-a_{x},b_{y}-a_{y}}
\end{align*}
Mit dem Satz des Pythagoras können die Längen der Vektoren und somit auch
deren Quadrate berechnet werden:
\begin{align*}
  |\vv{a}|^{2} &= a_{x}^{2}+a_{y}^{2} &
  |\vv{b}|^{2} &= b_{x}^{2}+b_{y}^{2}
\end{align*}
Analog kann die Länge der Seite $c$ aus den Komponenten von $\vv{a}$ und
$\vv{b}$ berechnet werden:
\[\begin{eqt}
  |\vv{c}|^{2} &= (b_{x}-a_{x})^{2}+(b_{y}-a_{y})^{2} & \text{ausmultiplizieren} \\[2mm]
  |\vv{c}|^{2} &= b_{x}^{2}-2a_{x}b_{x}+a_{x}^{2}+b_{y}^{2}-2a_{y}b_{y}+a_{y}^{2} & \text{Terme umstellen} \\[2mm]
  |\vv{c}|^{2} &= a_{x}^{2}+a_y^{2}+b_{x}^{2}+b_{y}^{2}-2a_{x}b_{x}-2a_{y}b_{y} & \text{ersetzen} \\[2mm]
  |\vv{c}|^{2} &= |\vv{a}|^{2}+|\vv{b}|^{2}-2a_{x}b_{x}-2a_{y}b_{y} & -2 \text{ ausklammern} \\[2mm]
  |\vv{c}|^{2} &= |\vv{a}|^{2}+|\vv{b}|^{2}-2\left(a_{x}b_{x}+a_{y}b_{y}\right)
\end{eqt}\]
Nach einigen Umformungen entsteht eine Gleichung, welche grosse Ähnlichkeit
mit dem Kosinussatz hat:
\[
  |\vv{c}|^{2} = |\vv{a}|^{2}+|\vv{b}|^{2}-2\cdot |\vv{a}|\cdot |\vv{b}|\cdot\cos(\gamma)
\]
Daraus kann geschlossen werden, dass folgendes gilt:
\[
  a_{x}b_{x}+a_{y}b_{y} = |\vv{a}|\cdot |\vv{b}|\cdot\cos(\gamma)
\]
Die linke Seite dieser Gleichung kann ausschliesslich aus den Komponenten der
Vektoren $\vv{a}$ und $\vv{b}$ berechnet werden.
\begin{theorem}
  \textbf{Definition Skalarprodukt 2D:} Das Skalarprodukt der Vektoren
  $\vv{a} = \vxy{a_{x},a_{y}}$ und $\vv{b} = \vxy{b_{x},b_{y}}$ ist
  wie folgt definiert:
  \[
    \vv{a}\circ \vv{b} = a_{x}\cdot b_{x}+a_{y}\cdot b_{y}
  \]
\end{theorem}
\begin{example}
  Das Skalarprodukt der beiden Vektoren $\vv{a} = \vxy{2,-3}$
  und $\vv{b} = \vxy{5,4}$ beträgt:
  \[
    \vv{a}\circ \vv{b}
    = \vxy{2,-3}\circ \vxy{5,4}
    = 2\cdot 5+(-3)\cdot 4
    = 10-12
    = -2
  \]
\end{example}
Analog wird das Skalarprodukt für dreidimensionale Vektoren definiert:
\begin{theorem}
  Das Skalarprodukt der Vektoren
  $\vv{a} = \vxyz{a_{x},a_{y},a_{z}}$ und $\vv{b} = \vxyz{b_{x},b_{y},b_{z}}$ ist
  wie folgt definiert:
  \[
    \vv{a}\circ \vv{b} = a_{x}\cdot b_{x}+a_{y}\cdot b_{y}+a_{z}\cdot b_{z}
  \]
\end{theorem}
\begin{example}
  Das Skalarprodukt der beiden Vektoren $\vv{a} = \vxyz{2,-3,1}$
  und $\vv{b} = \vxyz{-1,5,4}$ beträgt:
  \[
    \vv{a}\circ \vv{b}
    = \vxyz{2,-3,1}\circ \vxyz{-1,5,4}
    = 2\cdot (-1)+(-3)\cdot 5+1\cdot 4
    = -2-15+4
    = -13
  \]
\end{example}
Aus dem Vergleich mit dem Kosinussatz ergibt sich folgende Beziehung zwischen
dem Skalarprodukt und dem Winkel der Vektoren:
\begin{theorem}
  \textbf{Skalarprodukt und Winkel:} Für das Skalarprodukt und den von den
  Vektoren eingeschlossenen Winkel $\varphi$ gilt folgende Beziehung:
  \[
    \vv{a}\circ \vv{b} = |\vv{a}|\cdot |\vv{b}|\cdot \cos(\varphi)
  \]
\end{theorem}
Mit Hilfe dieser Gleichung kann der Winkel zwischen zwei Vektoren berechnet
werden.

\newpage
% ------------------------------------------------------------------------------
\subsection{Winkel zwischen zwei Vektoren}
Wird die obenstehende Gleichung umgeformt, so kann damit der Winkel zwischen
zwei Vektoren berechnet werden.
\begin{theorem}
  \textbf{Winkel zwischen zwei Vektoren:} Der Winkel $\varphi$ zwischen den
  Vektoren $\vv{a}$ und $\vv{b}$ wird
  wie folgt berechnet:
  \begin{align*}
    \cos(\varphi) &= \frac{\vv{a}\circ \vv{b}}{|\vv{a}|\cdot |\vv{b}|} &
    \varphi &= \arccos\left(\frac{\vv{a}\circ \vv{b}}{|\vv{a}|\cdot |\vv{b}|}\right)
  \end{align*}
  Insbesondere gilt:
  \begin{itemize}
    \item $\vv{a}$ und $\vv{b}$ stehen genau dann senkrecht aufeinander, wenn $\vv{a}\circ \vv{b} = 0$.
    \item $\vv{a}$ und $\vv{b}$ sind genau dann parallel zueinander, wenn $\vv{a}\circ \vv{b} = 1$.
  \end{itemize}
\end{theorem}

\begin{example}
  Der Winkel zwischen den Vektoren $\vv{a} = \vxy{-1,4}$ und $\vv{b} = \vxy{5,-1}$ soll berechnet werden.

  Zunächst wird das Skalarprodukt der beiden Vektoren berechnet:
  \[
    \vv{a}\circ \vv{b} = -1\cdot 5+4\cdot (-1) = -5-4 = -9
  \]
  Wenn das Skalarprodukt 0 oder 1 ergibt, dann stehen die Vektoren senkrecht beziehungsweise
  parallel zueinander. Der Winkel wäre somit \ang{90} bzw. \ang{0}. Ansonsten werden als
  nächstes die Längen der Vektoren ermittelt:
  \begin{align*}
    |\vv{a}| &= \sqrt{(-1)^{2}+4^{2}} = \sqrt{1+16} = \sqrt{17} \\
    |\vv{b}| &= \sqrt{5^{2}+(-1)^{2}} = \sqrt{25+1} = \sqrt{26}
  \end{align*}
  Nun kann der Winkel berechnet werden:
  \[
    \varphi = \arccos\left(\frac{\vv{a}\circ \vv{b}}{|\vv{a}|\cdot |\vv{b}|}\right)
    = \arccos\left(\frac{-9}{\sqrt{17}\cdot \sqrt{26}}\right) \approx \arccos\left(-0.428\right) \approx \ang{115.34}
  \]
\end{example}

